\documentclass[12pt]{article}
\usepackage{amsmath,amssymb}

\begin{document}
\section*{{2024 AMC 12A Problems/Problem 18}}
Source: AoPS Wiki

\subsection*{Solution}
Let the midpoint of $AC$ be $P$ . We see that no matter how many moves we do, $P$ stays where it is. Now we can find the angle of rotation ( $\angle APB$ ) per move with the following steps: \[AP^2=(\frac{1}{2})^2+(1+\frac{\sqrt{3}}{2})^2=2+\sqrt{3}\] \[1^2=AP^2+AP^2-2(AP)(AP)\cos\angle APB\] \[1=2(2+\sqrt{3})(1-\cos\angle APB)\] \[\cos\angle APB=\frac{3+2\sqrt{3}}{4+2\sqrt{3}}\] \[\cos\angle APB=\frac{3+2\sqrt{3}}{4+2\sqrt{3}}\cdot\frac{4-2\sqrt{3}}{4-2\sqrt{3}}\] \[\cos\angle APB=\frac{2\sqrt{3}}{4}=\frac{\sqrt{3}}{2}\] \[\angle APB=30^\circ\] Since Vertex $C$ is the closest one and \[\angle BPC=360-180-30=150\] Vertex C will land on Vertex B when $\frac{150}{30}+1=\fbox{(A) 6}$ cards are placed. (someone insert diagram maybe) ~lptoggled, minor Latex edits by eevee9406

\end{document}
